% Configuration optimisée pour une présentation professionnelle de mémoire de soutenance
% Cette approche est compatible avec pdfLaTeX ou XeLaTeX

% Définition des couleurs de base
\definecolor{titleColor}{RGB}{0, 51, 102}  % Bleu foncé pour les titres
\definecolor{subtitleColor}{RGB}{70, 130, 180}  % Bleu acier pour sous-titres
\definecolor{accentColor}{RGB}{0, 102, 204}  % Bleu vif pour les accents
\definecolor{lightAccent}{RGB}{235, 245, 255}  % Bleu très clair pour les fonds

% Packages nécessaires pour la mise en forme avancée
\usepackage{titlesec}    % Pour les titres
\usepackage{tocloft}     % Pour la table des matières
\usepackage{setspace}    % Pour l'interligne
\usepackage{enumitem}    % Pour les listes

% Style pour les chapitres - présentation élégante avec ligne horizontale
\titleformat{\chapter}[display]
  {\normalfont\huge\bfseries\color{titleColor}}
  {\flushright\normalsize\color{subtitleColor}\chaptertitlename\ \thechapter}
  {20pt}
  {\Huge\filcenter}
  [\vspace{0.5cm}\titlerule]

% Style pour les sections - avec ligne fine à gauche pour une meilleure lisibilité
\titleformat{\section}
  {\normalfont\Large\bfseries\color{titleColor}}
  {\thesection}
  {1em}
  {}
  [\vspace{0.3cm}\hrule height 0.5pt\vspace{0.3cm}]

% Style pour les sous-sections - avec espacement et mise en valeur
\titleformat{\subsection}
  {\normalfont\large\bfseries\color{subtitleColor}}
  {\thesubsection}
  {1em}
  {}
  [\vspace{0.2cm}]

% Style pour les sous-sous-sections - avec indentation et style distinctif
\titleformat{\subsubsection}
  {\normalfont\normalsize\bfseries\itshape\color{subtitleColor}}
  {\thesubsubsection}
  {1em}
  {}

% Espacement entre les titres et le texte
\titlespacing*{\chapter}{0pt}{50pt}{40pt}
\titlespacing*{\section}{0pt}{3.5ex plus 1ex minus .2ex}{2.3ex plus .2ex}
\titlespacing*{\subsection}{0pt}{3.25ex plus 1ex minus .2ex}{1.5ex plus .2ex}
\titlespacing*{\subsubsection}{0pt}{3.25ex plus 1ex minus .2ex}{1.5ex plus .2ex}

% Configuration des styles de paragraphe et de texte
\usepackage{microtype}  % Amélioration fine de la typographie
\usepackage{ragged2e}   % Meilleur alignement du texte
\usepackage{parskip}    % Espace entre paragraphes
\setlength{\parindent}{1.5em}  % Indentation des paragraphes
\setlength{\parskip}{0.8em}    % Espace entre paragraphes

% Configuration des légendes pour une présentation professionnelle
\usepackage{caption}
\captionsetup{
  font=small,
  labelfont={bf,color=titleColor},
  textfont=it,
  width=0.9\textwidth,
  format=hang,
  margin=10pt
}

% Configuration des en-têtes et pieds de page - style élégant et professionnel
\usepackage{fancyhdr}
\pagestyle{fancy}
\fancyhf{}
\fancyhead[L]{\small\textit{\nouppercase{\leftmark}}}
\fancyhead[R]{\small\thepage}
\fancyfoot[C]{\small\textsc{Mémoire de soutenance}}
\renewcommand{\headrulewidth}{0.4pt}
\renewcommand{\footrulewidth}{0.2pt}

% Style spécial pour la première page de chapitre
\fancypagestyle{plain}{
  \fancyhf{}
  \renewcommand{\headrulewidth}{0pt}
  \fancyfoot[C]{\small\thepage}
}

% Configuration pour les liens hypertexte - couleurs professionnelles
\usepackage{hyperref}
\hypersetup{
  colorlinks=true,
  linkcolor=titleColor,
  citecolor=accentColor,
  urlcolor=accentColor,
  bookmarksnumbered=true,
  bookmarksopen=true
}

% Styles personnalisés pour les environnements - présentation améliorée
\usepackage{mdframed}
\usepackage{tcolorbox}
\tcbuselibrary{skins,breakable}

% Environnement résumé élégant avec cadre léger
\newenvironment{resume}
  {\begin{tcolorbox}[
    enhanced,
    colback=lightAccent,
    colframe=accentColor,
    arc=0pt,
    boxrule=0.5pt,
    fonttitle=\bfseries,
    toptitle=3pt,
    bottomtitle=3pt,
    top=5pt,
    bottom=5pt,
    title=Résumé
  ]\small\itshape}
  {\end{tcolorbox}}

% Environnement pour encadrer un texte important
\newenvironment{encadre}
  {\begin{mdframed}[
    linewidth=1pt,
    linecolor=subtitleColor,
    backgroundcolor=white,
    innertopmargin=10pt,
    innerbottommargin=10pt,
    innerleftmargin=10pt,
    innerrightmargin=10pt
  ]}
  {\end{mdframed}}
  
% Environnement pour les définitions
\newenvironment{definition}[1]
  {\begin{tcolorbox}[
    enhanced,
    breakable,
    colback=white,
    colframe=titleColor,
    colbacktitle=titleColor,
    coltitle=white,
    fonttitle=\bfseries,
    attach boxed title to top left={yshift=-2mm, xshift=5mm},
    boxed title style={sharp corners},
    title=#1
  ]}
  {\end{tcolorbox}}
  
% Style amélioré pour les listes
\usepackage{enumitem}
\setlist{noitemsep, topsep=5pt, parsep=2pt, partopsep=0pt}
\setlist[itemize]{leftmargin=1.5em}
\setlist[enumerate]{leftmargin=1.5em}
  {\begin{center}\begin{tabular}{|p{0.95\textwidth}|}\hline\small}
  {\\\hline\end{tabular}\end{center}}

% Style amélioré pour les figures
\usepackage{float}
\floatstyle{boxed}
\restylefloat{figure}
